\begin{center}\begin{tikzpicture}[font=\small,>=stealth,line join=round]
\draw[gray!40,rounded corners=3pt] (0,0) rectangle (14.2,8.3);\node at (7.1,7.700000000000001) {两条轴分别配对，两次反射得到水平平移6};
\begin{scope}[shift={(1.9181818181818182,3.7409090909090907)},x=0.5454545454545454cm,y=1.3939393939393938cm]\draw[->,gray] (-1.5,0)--(9.5,0) node[right]{$x$};\draw[->,gray] (0,-1.5)--(0,1.8) node[above]{$y$};\draw[blue,thick,solid,domain=-1:9,samples=180] plot(\x,{cos(deg(pi*(\x-1)/3))});\fill[orange!85!black] (0,0.5) circle(2.2pt);\fill[orange!85!black] (2,0.5) circle(2.2pt);\fill[orange!85!black] (6,0.5) circle(2.2pt);\foreach \a in {1,4}{\draw[violet,dashed] (\a,-1.4)--(\a,1.4);\node[above,violet] at (\a,1.4){$x=\a$};}\end{scope}\node[anchor=west,align=left,text width=5.0cm] at (8.4,4.0) {示例 $f=\cos\frac{\pi(x-1)}3$\\[9pt]轴 $x=1,\ x=4$\\[9pt]$0\to2\to6$\\[9pt]高度都为 $\frac12$\\[9pt]周期 $T=2(4-1)=6$};
\node[text width=12.6cm,align=center] at (7.1,.65) {只讨论完整图像的两条轴；得出的6是一个周期，不一定最小。};
\end{tikzpicture}\end{center}
